\clearpage\SourceLeaf{196}%
\pdfbookmark[1]{Gauss to Schumacher, 1 April 1843}{schumacher18430401}
\fancyhead[L]{\small GAUSS TO SCHUMACHER}
\begin{center}
{\small GAUSS TO SCHUMACHER.\par}
\medskip\hrule\medskip
{\small Astronomische Nachrichten edited by H.\ C.\ \textsc{Schumacher}.\par
Vol.\ 20.\ No.\ 474.\ p.\ 299.\par}
\medskip\hrule
\end{center}\medskip
\begin{flushright}Göttingen, 1 April 1843.\end{flushright}
To determine a position at a given time from orbital elements, I prefer to use \textsc{Burckhardt}’s table to calculate the anomaly. It goes only as far as $163^0\ 45'$ and therefore no longer suffices for the comet’s present position according to Mr.\ \textsc{Galle}’s elements. \textsc{Barker}’s table covers the whole range, but the laborious interpolation makes it very inconvenient for large anomalies. In such cases I usually use a particular procedure that you might be glad to hear about. If $M$ is the number with which one would enter \textsc{Barker}’s table (or, for larger values, with their logarithms), so that $M=\frac{\text{elapsed time}}{nq^{\frac32}}$ where $\log n=0{,}0398723$, I set $\log\frac{MM}{16875}=3P$ and try, using my small logarithm table and taking $A$ and $B$ in the sense used there, to satisfy the equation $3A+2B=3P$. This can always be done very quickly when $P$ is large. If $a$ is then the number corresponding to the logarithm $A$, setting the anomaly $=v$ gives
\[
\operatorname{tang}\tfrac12v=\BareRoot{(3a)}\quad\text{or}\quad
\log\operatorname{tang}\tfrac12v=\tfrac12(A+\log3)
\]

The logarithm of the radius vector can then. also be calculated very easily. One goes back into the first column with $A+\log3$, or $A+\log3=A^*$, and the corresponding quantity in the second column $=B^*$ follows, immediately giving the logarithm of the radius vector $=A^*+B^*+\log q$.

\SourceLeaf{197}%
At least for the first trials, it is somewhat easier to solve that equation indirectly, almost by inspection, in the form $C=P+\frac13B$. One can first look up $P$ in the third column, or set $P=C'$ and set the corresponding quantity in the second column $=B'$, then take $P+\frac13B'=C''$ and the corresponding quantity from the second column of the table $=B''$, and then (where needed) take $P+\frac13B''=C'''$ and the corresponding $B'''$, and so on. The calculation comes to a halt very quickly. If the accuracy of five-figure logarithms is insufficient, \textsc{Matthiessen}’s table can be used to advantage here (I generally dislike using it because of the misplaced economy with which it was so needlessly printed). But I prefer to use it only after I have used the smaller table to find the two entries between which the final value of $A$ falls. I then prefer to use the equation in its original form $3A+2B=3P$.

\Needspace{9\baselineskip}
If, for example, the anomaly is to be determined for February $48.33333$, or for the time since perihelion $20^{\mathrm T}87663$, then according to \textsc{Galle}’s elements [Astr.\ Nachrichten vol.\ 20.\ No.\ 474 Berlin 25 March 1843]
\begin{center}
{\small\setlength{\tabcolsep}{3pt}
\begin{tabular}{r@{\ }r@{\qquad}r@{\ }r}
$q^{\frac32}\ .\ .\ .\ .\ .\ .$&$7.0809490$&$20.87663\ .\ .\ .\ .$&$1.3196604$\\
$n\BareRoot{16875}\ .\ .$&$2.1534942$&&$9.2344432$\\\cline{1-2}\cline{4-4}
Const.\ logarithm $=$&$9.2344432$&&$2.0852172$\\[3pt]
&&Thus $3P=$&$4.1704344$\\
&&$P=$&$1.3901448$
\end{tabular}}
\end{center}
The small tables give
\[
\begin{array}{ll}
B'=0.01806&\quad C''=1.39616\\
B''=0.01781&\quad C'''=1.39608
\end{array}
\]
which already settles the calculation and gives $A=1.37827$. \textsc{Matthiessen}’s table gives the more accurate value $A=1.3782739$. The calculation then proceeds as follows
\begin{center}
\begin{tabular}{r@{\ }l}
$A=$&$1.3782739$\\
$3\ .\ .\ .\ .\ .$&$0.4771213$\\\cline{2-2}
&$1.8553952$\\[3pt]
&$0.9276976=\log\operatorname{tang}\ 83^0\ 15'\ 49''53$\\
\multicolumn{2}{r}{and the true anomaly $=\qquad166\ 31\ 39{,}06$}
\end{tabular}
\end{center}

\SourceLeaf{198}%
Furthermore, the corresponding values are
\begin{center}
\begin{tabular}{r@{\ }r}
$A^*=$&$1.8553952$\\
$B^*=$&$0.0060170$\\
$q\ .\ .\ .\ .$&$8.0539660$\\\cline{1-2}
Logarithm of the radius vector $=$&$9.9153782$
\end{tabular}
\end{center}

One can also see that this method is simply an indirect solution of the familiar cubic equation relating the tangent of half the anomaly to the area of the sector. At the same time, one can see that my logarithm table, or \textsc{Matthiessen}’s for more precise calculations, can be used in much the same way to find very easily all the real roots of \textit{any algebraic equation with at most three effective terms}, as I recently showed for the quadratic equation in the latest edition of \textsc{Vega}’s logarithm table.
\par\medskip\centerline{\rule{.25\textwidth}{.3pt}}
